
\section{RQ1: Does single-transform adaptation compose?}

Our primary claim is that only composition present in the training data works, and it only works within the bounds of what that data covered. 

When attempting composition at the inference level via a learned router over per-transform specialists, the mixture yields zero benefit over a random gate (+0.0000). At the weight level, merging strategies such as TIES and DARE land at or below the clean-code-only control. However, at the data level, training breadth generalises to stacked inputs (+3.49, +2.90, and +3.05 points across three model scales). 

\begin{table}[ht]
\centering
\caption{Performance of composition strategies on stacked inputs.}
\begin{tabular}{@{}llp{8cm}@{}}
\toprule
\textbf{Level} & \textbf{Mechanism} & \textbf{Result} \\
\midrule
Inference & Learned router & +0.0000 [$-0.0081$, $+0.0081$] vs random gate. \\
Weights & TIES / DARE merges & Below clean-code control ($-3.13$ [$-4.78$, $-1.40$]). \\
Data & Breadth (\texttt{mono\_all}) & +3.49 / +2.90 / +3.05 on depth-2 composites; $-3.79$ / $-4.28$ / $-2.64$ on the unseen family.\\
\bottomrule
\end{tabular}
\end{table}

\begin{table}[ht]
\centering
\caption{Master Summary of Interventions. Performance is measured as the point difference ($\Delta$) compared to the clean-code-tuned baseline (\texttt{tuned\_L0}). The final column summarizes the percent of original (clean) code competence regained on obfuscated inputs without sacrificing generalization.}
\label{tab:master_summary}
\begin{tabular}{@{}lcccc@{}}
\toprule
\textbf{Intervention} & \textbf{Clean Code} & \textbf{Stacked Seen} & \textbf{Unseen Family} & \textbf{\% of Clean Regained} \\
\midrule
\textbf{Baseline (Clean Code)} & Reference & Reference & Reference & 100\% (Reference) \\
\midrule
\textbf{RQ1: Simple Approaches} & & & & \\
\quad Inference Routing & -- & $+0.00$ & -- & 0\% (vs random gate) \\
\quad Task-Vector Merging & -- & $-3.13$ & -- & Negative (below baseline) \\
\quad Training Breadth & $-2.60$ & $+3.49$ & $-3.79$ & Gains stacked; loses unseen \\
\midrule
\textbf{RQ2: Loss Objective} & & & & \\
\quad Paired-Consistency ($\lambda=3$) & $-0.30$ & $+4.76$ & $-0.30$ & Recovers stacked at zero cost \\
\bottomrule
\end{tabular}
\end{table}

Crucially, this breadth creates a dissociation: its gain on seen composites saturates by a quarter of the training corpus, while its tax on unseen families grows monotonically with data volume (scaling from $-0.49$ to $-3.79$ at full corpus). Providing more data buys nothing on covered domains while costing more on unseen ones. It is important to clarify that this is not a blanket failure of composition, but rather a rigid boundary on where data-driven composition succeeds.
